% CCSS 7.NS.A.3 — Solving Real-World Problems with Rational Numbers (AK: culturally responsive contexts)
\section{Solving Real-World Problems with Rational Numbers}

\begin{learningGoals}
\begin{itemize}
    \item Solve multi-step problems using all four operations with rational numbers
    \item Estimate first, solve, and check that answers make sense
\end{itemize}
\end{learningGoals}

\begin{conceptBox}{Putting It All Together}
Real-world problems often mix addition, subtraction, multiplication, and division — sometimes with fractions \textbf{and} decimals in the same problem.
\begin{itemize}[leftmargin=*]
    \item \textbf{Read carefully.} Decide which operation each part needs.
    \item \textbf{Watch the signs.} Gains and increases are positive. Losses and decreases are negative.
    \item \textbf{Estimate first.} Use rounded numbers to predict a reasonable answer.
    \item \textbf{Check your answer.} Does the sign make sense? Is the size reasonable?
\end{itemize}
\end{conceptBox}

\begin{workedExample}{Alaska Winter Temperatures}
At midnight in Fairbanks, the temperature is $-28.5°$F. By noon, it has risen $\frac{3}{4}$ of the way toward $0°$F. What is the noon temperature?

\textbf{Estimate:} $\frac{3}{4}$ of $28.5 \approx 21$. So about $-28.5 + 21 \approx -7°$F.

\textbf{Solve:} Rise $= \frac{3}{4} \times 28.5 = \frac{3}{4} \times \frac{57}{2} = \frac{171}{8} = 21.375$

Noon temp $= -28.5 + 21.375 = -7.125°$F
\answerHighlight{$-7.125°$F (about $-7°$F — estimate checks out)}
\end{workedExample}

\begin{realWorld}[Life in Alaska]
Alaskans use rational numbers every day: tracking temperature swings from $-40°$F to $60°$F, measuring elevation changes from sea level to mountain peaks, recording salmon catch data in fisheries, and managing budgets in rural communities. Math helps make sense of this beautiful, challenging environment.
\end{realWorld}

\mascotSays{When a problem has lots of steps, write down each step on its own line. It's much easier to catch sign errors!}

\begin{practiceBox}{Solving Real-World Problems}
\resetProblems

\prob A fishing boat catches $345.5$ lb of salmon on Monday and $278.75$ lb on Tuesday. The crew keeps $\frac{2}{3}$ of the total and sells the rest at $\$4.50$/lb. How much do they earn from the sale?
\answerExplain{$\$937.88$}{Total catch: $345.5 + 278.75 = 624.25$ lb. Sold: $624.25 \times \frac{1}{3} = 208.08\overline{3}$ lb. Revenue: $208.08\overline{3} \times 4.50 \approx \$937.88$.}

\prob The temperature in Barrow is $-18°$F at $6$~AM. It drops $5.5°$F by midnight and then rises $12.25°$F by noon the next day. What is the noon temperature?
\answerExplain{$-11.25°$F}{Start: $-18$. Drop: $-18 + (-5.5) = -23.5$. Rise: $-23.5 + 12.25 = -11.25°$F.}

\prob A hiker descends $1{,}250$ ft from a mountain summit at $4{,}380$ ft, rests, then climbs $\frac{1}{2}$ of the distance back up. What is the final elevation?
\answerExplain{$3{,}755$ ft}{After descent: $4{,}380 - 1{,}250 = 3{,}130$ ft. Climbs back $\frac{1}{2} \times 1{,}250 = 625$ ft. Final: $3{,}130 + 625 = 3{,}755$ ft.}

\prob A family budgets $\$1{,}200$ per month. They spend $\frac{2}{5}$ on housing, $\frac{1}{4}$ on food, and $\$180$ on fuel. How much is left?
\answerExplain{$\$240$}{Housing: $\frac{2}{5} \times 1{,}200 = \$480$. Food: $\frac{1}{4} \times 1{,}200 = \$300$. Spent: $480 + 300 + 180 = \$960$. Left: $1{,}200 - 960 = \$240$.}

\prob Evaluate: $-\dfrac{1}{2} \times 8 + \dfrac{3}{4} \div \left(-\dfrac{1}{2}\right)$
\answerExplain{$-5\frac{1}{2}$}{Multiply first: $-\frac{1}{2} \times 8 = -4$. Divide: $\frac{3}{4} \div (-\frac{1}{2}) = \frac{3}{4} \times (-2) = -\frac{3}{2}$. Add: $-4 + (-\frac{3}{2}) = -\frac{11}{2} = -5\frac{1}{2}$.}

\end{practiceBox}
